\vspace{3pt plus 2pt}\noindent\hspace*{1.2em}\kbvlabel{ಉದಾ}\par\vspace{1pt}

\begin{kbvderiv}
\vspace{4pt plus 3pt minus 1pt}
\begin{center}\begin{adjustbox}{max width=\textwidth}
{\renewcommand{\arraystretch}{1.22}\setlength{\tabcolsep}{4pt}
\begin{tabular}{l c l}
ಮಗು + ಕಳು & {\color{kbvgrey}$=$} & \kbvresult{ಮಕ್ಕಳು}{\color{kbvgrey} (ಬ.ವ.)} \\
ಮಗ + ಕಳು & {\color{kbvgrey}$=$} & \kbvresult{ಮಕ್ಕಳು} \\
ಮಗಳು + ಕಳು & {\color{kbvgrey}$=$} & \kbvresult{ಮಕ್ಕಳು} \\
ದೊಡ್ಡದು + ವು & {\color{kbvgrey}$=$} & \kbvresult{ದೊಡ್ಡವು} \\
ಇದು + ವು & {\color{kbvgrey}$=$} & \kbvresult{ಇವು} \\
ನಾನು + ವು & {\color{kbvgrey}$=$} & \kbvresult{ನಾವು} \\
\end{tabular}}
\end{adjustbox}\end{center}
\vspace{4pt plus 3pt minus 1pt}
\end{kbvderiv}

\kbvexercise{ಅಭ್ಯಾಸ ಪತ್ರಿಕೆ ೧೦}

\vspace{2.5pt plus 3pt minus 0.5pt}\noindent\hspace*{1.2em}\kbvnum{(1)}\hspace{0.55em}ವಚನಗಳಲ್ಲಿ ಎಷ್ಟು ಪ್ರಕಾರಗಳಿವೆ? ಅವು ಯಾವುವು?\par

\vspace{2.5pt plus 3pt minus 0.5pt}\noindent\hspace*{1.2em}\kbvnum{(2)}\hspace{0.55em}ಏಕವಚನಕ್ಕೂ ಬಹುವಚನಕ್ಕೂ ಇರುವ ವ್ಯತ್ಯಾಸವೇನು?\par

\vspace{2.5pt plus 3pt minus 0.5pt}\noindent\hspace*{1.2em}\kbvnum{(3)}\hspace{0.55em}ಗಳು, ಅರು, ಅಂದಿರು - ಈ ಆಗಮಗಳು ಸೇರಿ ಆಗುವ ಬಹುವಚನ ಶಬ್ದಗಳಿಗೆ ಒಂದೊಂದು ಉದಾಹರಣೆಯನ್ನು ಕೊಡಿ.\par

\vspace{2.5pt plus 3pt minus 0.5pt}\noindent\hspace*{1.2em}\kbvnum{(4)}\hspace{0.55em}ಈ ಶಬ್ದಗಳ ಬಹುವಚನ ರೂಪಗಳನ್ನು ತಿಳಿಸಿ - ಕುದುರೆ, ಸೈನಿಕ, ಅಜ್ಜ, ರಾಣಿ, ಅದು, ಅವನು, ಹೆಂಡತಿ.\par

\vspace{2.5pt plus 3pt minus 0.5pt}\noindent\hspace*{1.2em}\kbvnum{(5)}\hspace{0.55em}ಈ ಶಬ್ದಗಳ ಏಕವಚನ ರೂಪಗಳನ್ನು ತಿಳಿಸಿ - ಮುದುಕರು, ತಮ್ಮಂದಿರು, ಪಕ್ಷಿಗಳು, ಕಳ್ಳರು, ಚಿಕ್ಕವು, ಇವರು.\par

\vspace{2.5pt plus 3pt minus 0.5pt}\noindent\hspace*{1.2em}\kbvnum{(6)}\hspace{0.55em}“ನಿಮ್ಮ ಅಣ್ಣಂದಿರು ರಾಮನ ಪಾತ್ರವನ್ನು ಚೆನ್ನಾಗಿ ಅಭಿನಯಿಸಿದರು” - ಈ ವಾಕ್ಯದಲ್ಲಿ ಅಣ್ಣಂದಿರು ಎಂಬ ಶಬ್ದವನ್ನು ಯಾವ ಅರ್ಥದಲ್ಲಿ ಬಳಸಲಾಗಿದೆ?\par

\vspace{2.5pt plus 3pt minus 0.5pt}\noindent\hspace*{1.2em}\kbvnum{(7)}\hspace{0.55em}ತಂದೆಯವರು, ಸುಬ್ಬಮ್ಮನವರು, ರಾಮಾಶಾಸ್ತ್ರಿಗಳು, ದಿವಾನರು - ಈ ನಾಮಪದಗಳಲ್ಲಿ ಕಂಡುಬರುವ ವಿಶೇಷವೇನು?\par

\vspace{2.5pt plus 3pt minus 0.5pt}\noindent\hspace*{1.2em}\kbvnum{(8)}\hspace{0.55em}ಈ ಬಹುವಚನ ಶಬ್ದಗಳಿಂದ ಆಗಮಗಳನ್ನು ಪ್ರತ್ಯೇಕಿಸಿ ಕ್ರಮವಾಗಿ ಬರೆಯಿರಿ - ಮಾವಂದಿರು, ಹಾವಾಡಿಗರು, ಇಲಿಗಳು, ಜವಾನರು.\par
